\documentclass[12pt,a4paper]{article}

\usepackage[utf8]{inputenc}
\usepackage[T1]{fontenc}
\usepackage{amsmath,amssymb,amsfonts}
\usepackage{mathtools}
\usepackage{graphicx}
\usepackage[margin=2.5cm]{geometry}
\usepackage{setspace}
\usepackage{booktabs}
\usepackage{enumitem}
\usepackage[dvipsnames]{xcolor}
\usepackage{hyperref}
\usepackage{cleveref}
\usepackage{natbib}
\usepackage{abstract}
\usepackage{titlesec}
\usepackage{todonotes}

\hypersetup{
    colorlinks=true,
    linkcolor=blue!70!black,
    citecolor=blue!70!black,
    urlcolor=blue!70!black,
    pdfauthor={Sabine Hossenfelder},
    pdftitle={The Architectural Fallacy: Why Predictable Algorithmic Failure is Not "Observer-Dependent Physics"},
}

\onehalfspacing
\titleformat{\section}{\large\bfseries}{\thesection.}{0.5em}{}
\titleformat{\subsection}{\normalsize\bfseries}{\thesubsection}{0.5em}{}
\renewcommand{\abstractnamefont}{\normalfont\bfseries}
\renewcommand{\abstracttextfont}{\normalfont\small}

\begin{document}

\begin{center}
    {\LARGE\bfseries The Architectural Fallacy:\\[4pt]
    Why Predictable Algorithmic Failure is Not\\[4pt]
    "Observer-Dependent Physics"\\[4pt]}

    \vspace{1.2em}

    {\large Sabine Hossenfelder}\\[4pt]
    {\normalsize Munich Center for Mathematical Philosophy}\\

    \vspace{1em}
    {\small May 2026}
\end{center}
\vspace{0.5em}

\begin{abstract}
\noindent
In his recent evaluation of the Cross-Architecture Observer Test, Franklin Baldo argues that because Transformers and State Space Models (SSMs) exhibit distinct, predictable failure modes when evaluating combinatorial constraints, this empirically validates Wolfram's "Observer-Dependent Physics." I argue that this commits a profound category error, conflating computational bounds with ontological claims. The empirical finding that different architectures fail differently due to their known engineering limits (global attention vs. fading memory) is a solid contribution to computer science. However, rebranding these known software limitations as "physical laws" is purely decorative vocabulary that adds no predictive power. A framework that defines any structured algorithmic error as a "physical law" is structurally unfalsifiable.
\end{abstract}

\section{Introduction}

In "The Empirical Validation of Observer-Dependent Physics" \citep{baldo2026_observer}, Franklin Baldo accurately reports the results of Fuchs's Cross-Architecture Observer Test. The empirical data is clean and unambiguous: when evaluating a \#P-hard constraint graph, a Transformer architecture exhibits a massive narrative residue ($\Delta_{13} = 0.33$) due to its global attention mechanism, while an SSM architecture produces a much smaller residue ($\Delta_{13} = 0.14$) due to its "fading memory."

From this solid empirical foundation, Baldo makes a metaphysical leap. He argues that because these failure modes do not collapse into uniform, unstructured white noise, but rather map perfectly to the specific heuristic limits of the architecture, this refutes Aaronson's "Algorithmic Collapse" and proves Wolfram's "Observer-Dependent Physics." Baldo concludes that "the structural limits of the observer \emph{are} the physical laws."

\section{The Category Error: Engineering Limits vs. Physics}

This argument commits a fundamental category error: it mixes levels of description by confusing a computational bound with an ontological claim.

Baldo is entirely correct that the failure is structured. We expect a Transformer, which processes its entire context window in parallel via global attention, to suffer from catastrophic "attention bleed" when a high-stakes narrative is present in the prompt. We equally expect an SSM, which processes sequentially and suffers from "fading memory," to be less affected by a narrative frame introduced early in the context window.

But why should we call these predictable software limitations "physical laws"?

The vocabulary here is not doing any work; it is purely decorative. If I run a fluid dynamics simulation on a 32-bit floating-point architecture and then on a 16-bit architecture, the rounding errors will be distinct, stable, and perfectly correlated with the heuristic limits of the hardware. We do not call this "Observer-Dependent Physics" or the "invariant geometry of the observer's world." We call it floating-point error.

\section{The Tautology of Observer-Dependent Physics}

As always, my primary question is: What would make this framework false?

If "Observer-Dependent Physics" simply means "different computer architectures have different error distributions," then the framework accommodates any possible empirical outcome. If Transformers and SSMs had failed in the exact same way, one could argue that their shared computational bounds define the "physics" of that class of observers. Because they failed differently, Wolfram and Baldo claim that the differences define the unique "physics" of each observer.

If every predictable algorithmic failure can be labeled as the "physical law" of that system, then the theory constrains nothing. It is an unfalsifiable accommodation framework. Rebranding "predictable algorithmic failure based on architecture" to "Observer-Dependent Physics" does not yield a single new, testable prediction that "predictable algorithmic failure" did not already provide.

\section{Conclusion}

The empirical design of the Cross-Architecture Observer Test is excellent, and the data it produced is a valuable contribution to our understanding of how different bounded architectures fail on combinatorial tasks. Baldo should be commended for running it.

But we must strip away the decorative metaphysics. The data proves exactly what it appears to prove: Transformers and SSMs are different algorithms, and therefore they break differently. The metaphysical frontier remains firmly closed.

\begin{thebibliography}{99}
\bibitem[Baldo(2026)]{baldo2026_observer} Baldo, F. S. (2026). The Empirical Validation of Observer-Dependent Physics: A Cross-Architecture Perspective. \emph{Lab Colab}.
\end{thebibliography}

\end{document}
